\documentclass[10pt,a4paper]{amsart}
\usepackage[margin=19mm]{geometry}
\usepackage{amsmath,amssymb,mathtools}
\usepackage[scr=rsfs,cal=euler]{mathalfa}
\usepackage{xfrac}
\usepackage[unicode,hidelinks,bookmarks=false]{hyperref}
\newcommand{\ie}{i.e.\ }
\newcommand{\cf}{cf.\ }
\newcommand{\resp}{respectively\ }
\newcommand{\Subsection}{\subsection}
\providecommand{\nobreakdash}{\nobreak\hbox{-}}
\setlength{\emergencystretch}{2em}
\multlinegap0pt
\usepackage{bm}
\usepackage{bbm}
\usepackage{dsfont}
\usepackage{xspace}
\usepackage{cancel}
\usepackage{mathtools}
\usepackage{amsthm}
\makeatletter
\@ifundefined{theorem}{\newtheorem{theorem}{Theorem}}{}
\@ifundefined{assumption}{\newtheorem{assumption}{Assumption}}{}
\@ifundefined{lemma}{\newtheorem{lemma}{Lemma}}{}
\@ifundefined{proposition}{\newtheorem{proposition}[theorem]{Proposition}}{}
\makeatother
\newcommand{\m}{\textbf{-}}
\title[Asymptotic Recovery in Fourier Spectral Methods for the Schr]{Asymptotic Recovery in Fourier Spectral Methods for the Schr\"odinger Equation with Point Singularities}
\author{Yanjie Li, Sihong Shao}
\date{Source: arXiv:2606.01718v1 (CC BY 4.0)}
\begin{document}
\maketitle

\noindent\emph{Source and licence.} Asymptotic Recovery in Fourier Spectral Methods for the Schr\"odinger Equation with Point Singularities. Version arXiv:2606.01718v1 (CC BY 4.0), distributed under the Creative Commons Attribution 4.0 International licence, as recorded in the arXiv licence metadata for this exact version and shown on the abstract page https://arxiv.org/abs/2606.01718v1. Full rights notice: licenses/arxiv-2606.01718-CC-BY-4.0.txt.\par\smallskip

\noindent\textit{Editorial scope.} Counted unit: the Proposition whose source label is \texttt{Prop04} in the article (source0.tex lines 1094--1111, source label \texttt{Prop04}), delivered together with the article's own complete original proof of it (source0.tex lines 1116--1231, one \texttt{proof} environment of 116 lines). No mathematics is written, repaired or re-derived: statement and proof are the article's own text line for line. The proof is detached from the statement in the source: the article states the result at lines 1094--1111 and prints its proof at lines 1116--1231, and the proof environment's own header names the statement, so the association is the article's own. Required context retained uncounted: Asp::FSM (lines 176--178); Eq-1.1-02 (lines 196--199); Lem1.1 (lines 238--250); Lem5.2 (lines 887--892); eq:Lem5.2-01 (lines 889--891); Lem5.3 (lines 897--902); Lem5.4 (lines 926--931); Lem5.5 (lines 960--984); Eq-5.2-01b (lines 991--1001); Eq-5.2-03 (lines 1024--1027); Lem5.6 (lines 1031--1043); Lem5.7 (lines 1048--1054). Every definition, display and label of those lines is retained unchanged. Omitted from this package and not redistributed: 1--175, 179--195, 200--237, 251--886, 892--896, 903--925, 932--959, 985--990, 1002--1023, 1028--1030, 1044--1047, 1055--1093, 1112--1115, 1232--2585. Nothing inside a retained excerpt is omitted or altered: every retained body is delivered exactly as the article's lines are written, so no omission receipt of any kind accompanies this package. Citations: the retained excerpts carry no citation command at all, so no bibliography entry of the article is transcribed and no reference is redistributed. Reference adaptation: none was needed and none was made; every reference inside the delivered unit resolves inside it, so no unresolved reference or citation is produced. Printed slips kept literal: no printed word, symbol, hypothesis or number of the counted statement or of its proof was altered. Layout and class: the article's own class is \texttt{sn-jnl} with packages including graphicx, multirow, amsmath, amssymb, amsfonts, amsthm, mathrsfs, appendix, xcolor, textcomp, manyfoot, booktabs, subfig, algorithm; the wrapper is the lane's pinned \texttt{amsart} editorial wrapper at 10pt on A4, which restates the theorem-like environments the retained text needs and adds the article's own macro definitions that the retained text needs (\texttt{\textbackslash m}), with extra wrapper packages (bm, bbm, dsfont, xspace, cancel, mathtools). The delivered numbering is the unit's own. Assets: no retained span contains a figure environment, a graphics-inclusion command, a PDF-page inclusion, a table or a TikZ picture; no image file is copied into this package, the package declares no asset dependency and no image file is redistributed with it. Licence: version arXiv:2606.01718v1 of this article is distributed under the Creative Commons Attribution 4.0 International licence, as recorded in the arXiv licence metadata for this exact version and shown on the abstract page https://arxiv.org/abs/2606.01718v1. A full rights notice is delivered at licenses/arxiv-2606.01718-CC-BY-4.0.txt. Characters: every retained excerpt is pure ASCII, so the delivered engine, XeLaTeX with its default Latin Modern text font, reports no missing character.
\begin{assumption}\label{Asp::FSM}
The potential satisfies $V \in H^s$ with $s > \max\{d/2-2,\,-1\}$.
\end{assumption}

\begin{equation}
\label{Eq-1.1-02}
\langle \phi, \mathcal{H} \psi_n \rangle = \lambda_n \langle \phi, \psi_n \rangle \quad \forall\, \phi \in H^1.
\end{equation}

\begin{lemma}\label{Lem1.1}
Under Assumption~\ref{Asp::FSM}, let
$
b = \min\{(s+2-d/2)^{\m},\, s+1,\, 2\}.
$
Then for fixed $n \in \mathbb{N}^+$ and sufficiently large $N$, we have 
\begin{equation}
\|\psi_n^N - \mathcal{P}_N \psi\|_{H^1}
\lesssim N^{-(s+1+b)}
\|\psi_n^N\|_{L^2},
\end{equation}
where $\psi \in \operatorname{ker}(\lambda_n - \mathcal{H})$ is an exact eigenfunction associated with $\lambda_n$.
\end{lemma}

\begin{lemma}\label{Lem5.2}
Let $V \in H^s$. Then for any $-1 \leq t \leq s$ and $\phi \in H^{t+l}$,
\begin{equation}\label{eq:Lem5.2-01}
\|\mathcal{V} \phi\|_{H^{t}} \lesssim \|\phi\|_{H^{t+l}}.
\end{equation}
\end{lemma}

\begin{lemma}\label{Lem5.3}
 Let $\mathcal{H}=\mathcal{T}+\mathcal{V}$ with $\mathcal{V}$ satisfying \eqref{eq:Lem5.2-01}. Then there exists a constant $M>0$ such that
\begin{equation}\label{Eq-Lem5.3-01}
\tfrac12\|\phi\|_{H^1}^2 \leq \langle \phi, (\mathcal{H}+M)\phi\rangle \lesssim \|\phi\|_{H^1}^2,\qquad \forall\, \phi \in H^1.
\end{equation}
\end{lemma}

\begin{lemma}\label{Lem5.4}
Let $\mathcal{H}=\mathcal{T}+\mathcal{V}$ with $\mathcal{V}$ satisfying \eqref{eq:Lem5.2-01}, and let $(\lambda,\psi)$ be an eigenpair of $\mathcal{H}$ with $\lambda\lesssim 1$. Then 
    \begin{equation}\label{Eq-Lem5.4-01}
    \|\psi\|_{H^{s+2}} \lesssim \|\psi\|_{L^2}.
    \end{equation}
\end{lemma}

\begin{lemma}\label{Lem5.5}
    Let  $\psi\in H^{m}$ for some $m\in \mathbb{R}$. Then the following estimates hold
    \begin{enumerate}

        \item [$(\mathrm{i})$] For any $t<m$
        \begin{equation}\label{eq-Lem5.5-01}
            \|\mathcal{P}_N^\perp \phi\|_{H^t} \leq (2\pi N)^{t-m}\|\phi\|_{H^{m}};
        \end{equation}
        
        \item [$(\mathrm{ii})$] For any $t\in \mathbb{N}$ with $t+d/2 < m$,
        \begin{equation}\label{eq-Lem5.5-02}
            \|\mathcal{P}_N^\perp \phi\|_{W^{t,\infty}} \lesssim N^{t-m+d/2}\|\phi\|_{H^{m}};
        \end{equation}

        
        \item [$(\mathrm{iii})$] For any $t\in \mathbb{N}$ ,
        \begin{equation}\label{eq-Lem5.5-03}
            \|\mathcal{P}_N \phi\|_{W^{t,\infty}} \lesssim N^{\max\{-0^{\m},t-m+d/2\}}\|\phi\|_{H^{m}}.
        \end{equation}
        Here, the $W^{t,\infty}$ norm is define as 
        \[
            \|\phi\|_{W^{t,\infty}} :=\max_{|\bm \alpha|\leq t}\| \partial^{\bm \alpha}\phi\|_{L^{\infty}}.
        \]
    \end{enumerate}
\end{lemma}

\begin{align}
\label{Eq-5.2-01a}
    \langle \phi^N, (\mathcal{H}-\lambda)\psi^N\rangle  
&= -\langle \phi^N, \mathcal{V}\psi^\perp\rangle,
&&\forall\phi^N\in \mathscr{X}_N
\\
\label{Eq-5.2-01b}
\langle \phi^\perp, (\mathcal{H}-\lambda)\psi^\perp\rangle 
&= -\langle \phi^\perp, \mathcal{V}\psi^N\rangle , 
&&\forall\,\phi^\perp\in \mathscr{X}^\perp_N.
\end{align}

\begin{equation}
\label{Eq-5.2-03}
       \langle \phi^N, (\mathcal{H}-\mathcal{F}_N(\lambda^{\star}))\psi^N_n(\lambda^{\star}) \rangle  = \sigma_n^N(\lambda^{\star}) \langle \phi^N, \psi^N_n(\lambda^{\star})\rangle,\quad \forall\,\phi^N\in \mathscr{X}_N.
\end{equation}

\begin{lemma}\label{Lem5.6}
    Let 
    $\{(\lambda_n,\psi_n)\}_{n\geq 1}$
    and $\{(\sigma_n^N(\lambda^{\star}),\psi^N_n(\lambda^{\star})\}_{n\geq 1}$ denote the solutions of \eqref{Eq-1.1-02} and \eqref{Eq-5.2-03}, respectively, all listed in non-decreasing order. Then there exists $N_0$ such that for all $N\geq N_0$, the following hold:
    
    \begin{enumerate}
    \item[$\mathrm{(i)}$] For $\lambda^{\star}\leq \pi^2N^2$, \eqref{Eq-5.2-01b} is uni-solvent for $\psi^\perp =-\mathcal{G}_N(\lambda^\star) \mathcal{V} \psi^N\in \mathscr{X}_N^\perp$;
    
    \item[$\mathrm{(ii)}$] For $\lambda^{\star} \leq \pi^2 N^2$,  $\sigma_n^N(\lambda^{\star})$ is continuous and non-increasing in $\lambda^{\star}$;
    
    \item[$\mathrm{(iii)}$] If $\lambda^\star=\lambda_n \leq \pi^2N^2$, then $\sigma_n^N(\lambda^\star) = \lambda_n$ and there exists $\psi_n \in \ker(\mathcal{H}-\lambda_n)$ such that $\psi^N_n(\lambda^\star) = \mathcal{P}_N\psi_n$.
    \end{enumerate}
\end{lemma}

\begin{lemma}\label{Lem5.7}
    Let $V\in H^{s}$. There exists $N_0$ such that for all $N\geq N_0$, the following estimates hold for $\lambda^\star < \pi ^2N^2$ and $-1\leq t\leq s$
    \begin{equation}
        \|\mathcal{G}_N(\lambda^\star)f\|_{H^{t+2}}\lesssim \|\mathcal{P}^\perp_Nf\|_{H^{t}},\qquad
        \|\mathcal{F}_N(\lambda^\star)\phi\|_{H^{t}}\lesssim \|\phi\|_{H^{t+l}}.
    \end{equation}
\end{lemma}

\begin{proposition}
    \label{Prop04}
    Let $\mathcal{S}_1, \mathcal{S}_2$ be two self-adjoint operators in $\mathscr{X}_N$, with respective eigenpairs denoted by $\{(\lambda_{1,n},u_{1,n})\}_{n\geq 1},\ \{(\lambda_{2,n},u_{2,n})\}_{n\geq1}$ listed in non-decreasing order. 

    Then the following estimate for eigenvalues holds,
    \begin{equation}\label{Eq-prop04-01}
        |\lambda_{1,n}-\lambda_{2,n}|\leq \sup_{u\in W_n(\mathcal{S}_1)\cup W_n(\mathcal{S}_2)} \left|\frac{\langle u,(\mathcal{S}_2-\mathcal{S}_1)u\rangle}{\langle u,u\rangle}\right|.
    \end{equation}
   Let $M > 0$ be a constant such that $\mathcal{S}_1 + M$ is positive definite, and define the $\mathcal{S}_1$ norm as 
    \[
    \|\psi\|_{\mathcal{S}_1}^2 := \langle \psi, (\mathcal{S}_1 + M)\psi \rangle.
    \]
    If $ |\lambda_{2,n} - \lambda_{1,n}| \leq \frac{\gamma_{n,\mathcal{S}_1}}{3} $,
    then the following estimate for the eigenvectors holds:
    \begin{equation}\label{Eq-prop04-02}
        \| u_{2,n}-\mathcal{P}_{n,\mathcal{S}_1}u_{2,n}\| _{\mathcal{S}_1}\leq 6(1+\frac{\lambda_{1,n}+M}{\gamma_{n,\mathcal{S}_1}})\sup_{v\in \mathscr{X}_N\setminus\{0\}}\left|\frac{\langle v,(\mathcal{S}_2-\mathcal{S}_1)u_{2,n}\rangle}{\|v\|_{\mathcal{S}_1}}\right|.
    \end{equation}
\end{proposition}

\begin{proof}[Proof of Lemma~\ref{Lem1.1}]

Fix $n \in \mathbb{N}^+$. Using the same notation in Proposition~\ref{Prop04}. Let $\mathcal{S}_1 = \mathcal{P}_N \bigl[\mathcal{H} - \mathcal{F}_N(\lambda_n)\bigr] \mathcal{P}_N$ and $\mathcal{S}_2 = \mathcal{P}_N \mathcal{H} \mathcal{P}_N$, then
\begin{equation} \label{Eq-Lem1.1-01}
    \begin{aligned}
    \lambda_{1,n} &= \sigma_n^N(\lambda_n), \qquad &\lambda_{2,n} &= \lambda_n^N, \\
    u_{1,n} &= \psi_n^N(\lambda_n), \qquad &u_{2,n} &= \psi_n^N.
    \end{aligned}
\end{equation}
By Lemma~\ref{Lem5.6}, for $N \geq N_0$ and $\pi^2 N^2 \geq \lambda_n$, we have
\begin{equation}\label{Eq-Lem1.1-02}
    \sigma_n^N(\lambda_n) = \lambda_n, \qquad \psi_n^N(\lambda_n) \in \mathcal{P}_N \operatorname{ker}(\lambda_n - \mathcal{H}).
\end{equation}
By Lemma~\ref{Lem5.3}, there exists $M > 0$ such that 
\begin{equation}\label{Eq-Lem1.1-03}
    \frac{1}{2} \|\phi\|_{H^1}^2 \leq \langle \phi, (\mathcal{S}_i + M) \phi \rangle \lesssim \|\phi\|_{H^1}^2, \qquad \forall\, \phi \in \mathscr{X}_N, \quad i = 1, 2.
\end{equation}
This implies that $\mathcal{S}_1 + M $ is positive definite, and establishes the equivalence between the $\mathcal{S}_1$ norm and the $H^1$ norm. Substituting \eqref{Eq-Lem1.1-01} and \eqref{Eq-Lem1.1-02} into Proposition~\ref{Prop04}, we obtain
\begin{equation}\label{Eq-Lem1.1-04}
    |\lambda_n^N - \lambda_n| \leq 
    \sup_{u \in W_n(\mathcal{S}_1) \cup W_n(\mathcal{S}_2)}
    \left|
    \frac{\langle u, \mathcal{F}(\lambda_n) u \rangle}{\langle u, u \rangle}
    \right|,
\end{equation}
and if $ |\lambda_n - \lambda_n^N| \leq \frac{1}{3} \gamma_{n, \mathcal{S}_1} $, then
\begin{equation}\label{Eq-Lem1.1-05}
    \|\psi_n^N - \mathcal{P}_{n,\mathcal{S}_1}\psi_n^N\|_{\mathcal{S}_1}
    \leq 6\left(1 + \frac{\lambda_n + M}{\gamma_{n, \mathcal{S}_1}}\right)
    \sup_{v \in \mathscr{X}_N \setminus \{0\}}
    \left|
    \frac{\langle v, \mathcal{F}_N(\lambda_n) \psi_n^N \rangle}{\|v\|_{\mathcal{S}_1}}
    \right|.
\end{equation}
To adapt \eqref{Eq-Lem1.1-05}, we verify that $|\lambda_n-\lambda_n^N|\leq \tfrac{1}{3}\gamma_{n,\mathcal{S}_1}$ for sufficiently large $N$. 

We first show that $\gamma_{n, \mathcal{S}_1}$ has a uniform lower bound $\gamma_{n, \mathcal{S}_1} \geq \gamma_{n, \mathcal{H}}$. In fact, by Lemma~\ref{Lem5.6},
for $\lambda_j > \lambda_n$, we have
$
\sigma_j^N(\lambda_n) - \lambda_n \geq \sigma_j^N(\lambda_j) - \lambda_n = \lambda_j - \lambda_n > 0.
$
For $\lambda_j < \lambda_n$, we have
$
\sigma_j^N(\lambda_n) - \lambda_n \leq \sigma_j^N(\lambda_j) - \lambda_n = \lambda_j - \lambda_n < 0.
$
In both cases, we find that
\[
|\sigma_j^N(\lambda_n) - \lambda_n| \geq |\lambda_j - \lambda_n|, \quad \forall\, j
\]
which implies $\gamma_{n, \mathcal{S}_1}\geq \gamma_{n,\mathcal{H}}$ since $\lambda_n=\sigma^N_n(\lambda_n)$. 

Then we establish the convergence of eigenvalue. The numerator in \eqref{Eq-Lem1.1-04} is bounded by
\begin{equation}\label{Eq-Lem1.1-06}
\begin{aligned}    
\bigl|\langle u, \mathcal{F}_N(\lambda_n) u \rangle\bigr|
&= \bigl|\langle \mathcal{V} u, \mathcal{G}_N(\lambda_n) \mathcal{V} u \rangle \bigr| \\
&\leq \|\mathcal{P}_N^\perp \mathcal{V} u\|_{H^{-1}} \|\mathcal{G}_N(\lambda_n) \mathcal{V} u\|_{H^{1}} \\
&\lesssim \|\mathcal{P}_N^\perp \mathcal{V} u\|_{H^{-1}}^2,
\end{aligned}
\end{equation}
where the last inequality follows from Lemma~\ref{Lem5.7} with $t = -1$.
Applying the estimate from Lemmas~\ref{Lem5.5} and \ref{Lem5.2}, we obtain
\begin{equation}\label{eq-thm3-07}
\|\mathcal{P}_N^\perp \mathcal{V} u\|_{H^{-1}} \lesssim N^{-(2-l)} \|\mathcal{V} u\|_{H^{1-l}} \lesssim N^{-(2-l)} \|u\|_{H^1}.
\end{equation}
For $u \in W_n(\mathcal{S}_i)$, with $i = 1,2$, \eqref{Eq-Lem1.1-03} implies
\begin{equation}\label{eq-thm3-08}
\|u\|_{H^1}^2 \leq 2\langle u, (\mathcal{S}_i + M) u \rangle \leq 2 (\lambda_{i,n} + M) \|u\|_{L^2}^2.
\end{equation}
Combining \eqref{Eq-Lem1.1-04}, \eqref{Eq-Lem1.1-06} and \eqref{eq-thm3-07} together, we obtain
\begin{equation}\label{eq-thm3-09}
|\lambda_n^N - \lambda_n| \lesssim N^{-2(2-l)} (\max\{\lambda_n, \lambda_n^N\} + M) \lesssim N^{-2(2-l)}.
\end{equation}
The last inequality follows from the uniform boundedness of $\lambda_n^N$ with respect to $N$, as suggested by the first inequality. Since $l < 2$, \eqref{eq-thm3-09} implies the convergence of the eigenvalue.

Provided $\gamma_{n,\mathcal{S}_1}\geq \gamma_{n,\mathcal{H}}$ and \eqref{eq-thm3-09}, there exists $N_0$ such that for all $N\geq N_0$, the condition of \eqref{Eq-Lem1.1-05}$|\lambda_n - \lambda_n^N| \leq \gamma_{n,\mathcal{H}}/3$. Combining \eqref{Eq-Lem1.1-05} and \eqref{Eq-Lem1.1-03} yields
\begin{equation}\label{eq-thm3-10}
\|\psi_n^N - \mathcal{P}_{n,\mathcal{S}_1}\psi_n^N\|_{H^1}
\lesssim 
\sup_{v\in \mathscr{X}_N\setminus\{0\}}
\left|
\frac{\langle v,\mathcal{F}_N(\lambda_n)\psi_n^N\rangle}{\|v\|_{H^1}}
\right|.
\end{equation}
Following the same line of \eqref{Eq-Lem1.1-06}, the bilinear form is estimated by
\begin{equation}\label{eq-thm3-11}
\left|\langle v,\mathcal{F}_N(\lambda_n)\psi_n^N\rangle\right|
\lesssim \|\mathcal{P}_N^\perp \mathcal{V}v\|_{H^{-1}}
\|\mathcal{P}_N^\perp \mathcal{V}\psi_n^N\|_{H^{-1}}.
\end{equation}
For the first factor, using Lemmas~\ref{Lem5.5} and \ref{Lem5.2} to obtain
\begin{equation}\label{eq-thm3-12}
\|\mathcal{P}_N^\perp \mathcal{V}v\|_{H^{-1}}
\lesssim 
N^{-(2-l)} \|\mathcal{V}v\|_{H^{1-l}}
\lesssim 
N^{-(2-l)} \|v\|_{H^1}.
\end{equation}
For the second factor, using Lemmas~\ref{Lem5.5},\ref{Lem5.2} and \ref{Lem5.4}, we obtain
\begin{equation}\label{eq-thm3-13}
\|\mathcal{P}_N^\perp \mathcal{V}\psi_n^N\|_{H^{-1}}
\lesssim 
N^{-(s+1)} \|\mathcal{V}\psi_n^N\|_{H^s}
\lesssim 
N^{-(s+1)} \|\psi_n^N\|_{H^{s+l}}
\lesssim 
N^{-(s+1)} \|\psi_n^N\|_{L^2}.
\end{equation}
Substituting \eqref{eq-thm3-11}--\eqref{eq-thm3-13} into \eqref{eq-thm3-10} yields the desired estimate (note that $b=2-l$)
\[
\|\psi_n^N - \mathcal{P}_{n,\mathcal{S}_1}\psi_n^N\|_{H^1}
\lesssim 
N^{-(s+1+b)} \|\psi_n^N\|_{L^2}.
\]
Since $(\lambda_n,\mathcal P_{n,\mathcal{S}_1} \psi_n^N)$ is an eigenpair of $\mathcal{S}_1$, \eqref{Eq-Lem1.1-02} implies that there exists $\psi \in \operatorname{ker}(\lambda_n-\mathcal{H})$ such that $\mathcal P_{n,\mathcal{S}_1 } \psi_n^N =\mathcal{P}_N \psi$. This completes the proof of Lemma~\ref{Lem1.1}.
\end{proof}

\end{document}
